\section{Konzept und Pipeline-Idee}

\subsection{Leitprinzip und Datenfluss}

Die Pipeline trennt die Bilddomäne, die semantische Domäne und die
Geometriedomäne. Der Ablauf ist:

\begin{enumerate}
  \item Aus dem Rohvideo werden Arbeitsframes erzeugt.
  \item SAM~3.1 propagiert den Objektprompt und speichert hierarchische
        Masken, das heißt die unveränderte \texttt{default}-Maske sowie die
        erodierten Ableitungen \texttt{middle} und \texttt{small}.
  \item COLMAP schätzt aus unmaskierten Frames ein ursprüngliches
        Kameramodell und eine sparse Punktwolke.
  \item Bei verzeichnungsbehafteten Kameramodellen wird eine ideale PINHOLE-
        Szene erzeugt; die Masken werden mit derselben Abbildung in diese
        Szene gewarpt, das heißt per Pixel-Neuzuordnung (Nearest Neighbor)
        auf die undistortierten Bilder übertragen, damit die binäre
        Maskensemantik erhalten bleibt.
  \item STS trainiert die objektspezifische Gaussian-Repräsentation.
  \item Die Hochopazitätskopie der Objekt-Gaussians wird über die gewählte
        Meshroute in ein Coarse-Mesh überführt.
  \item DGtal extrahiert daraus eine lokale Centerline; anschließend können
        B-Spline und GeoJSON erzeugt werden.
\end{enumerate}

Abbildung~\ref{fig:pa-pipeline-flow} fasst diesen Datenfluss zusammen: Die
beiden aktuellen Meshrouten A und \texttt{sugar\_coarse} werden explizit
getrennt dargestellt. Der CloudCompare-Schritt ist als optionaler
Georeferenzierungszweig enthalten. Für den Testdatensatz existieren
CloudCompare-Transformationsmatrizen; sie wurden jedoch ohne unabhängige
Vermessungsreferenz erzeugt und sind daher kein validierter
PA-Genauigkeitsnachweis. Die Referenzierung ist zunächst rein lokal, das
heißt ohne Verschiebung in ein globales Koordinatensystem definiert. Das ist
methodisch unproblematisch: Eine UTM-Referenz lässt sich nachträglich
herstellen, indem die lokalen Koordinaten um einen frei wählbaren
UTM-Ankerpunkt versetzt und die vorhandene 4x4-Transformationsmatrix
anschließend unverändert angewandt wird; Rotation, Maßstab und Form der
Geometrie bleiben dabei erhalten. Was damit nicht entsteht, ist eine
unabhängige Genauigkeitsaussage zur absoluten Lage -- diese bleibt der
Bachelorarbeit vorbehalten.

\begin{figure}[H]
\centering
\resizebox{\textwidth}{!}{%
\begin{tikzpicture}[
      node distance=1.35cm and 2.0cm,
      box/.style={rectangle, draw=blue!65, fill=blue!5, very thick,
            minimum width=4.4cm, minimum height=0.9cm, align=center,
            rounded corners, text width=4.2cm},
      route/.style={rectangle, draw=purple!70, fill=purple!7, very thick,
            minimum width=4.6cm, minimum height=0.9cm, align=center,
            rounded corners, text width=4.4cm},
      decision/.style={diamond, draw=red!65, fill=red!7, very thick,
            text width=2.7cm, align=center, aspect=2.0},
      arrow/.style={-{Stealth[scale=1.05]}, thick},
      container/.style={rectangle, draw, dashed, rounded corners, inner sep=0.35cm}
]
      \node[box] (input) {Rohvideo\\Arbeitsframes\\Prompt \texttt{pipe}};
      \node[box, below=of input] (sam) {\textbf{Container A: SAM 3.1}\\Temporale Segmentierung\\default / middle / small};
      \node[box, below=of sam] (colmap) {\textbf{Container B: COLMAP}\\Unmaskiertes SfM\\Originales Kameramodell und Sparse-Modell};
      \node[box, below left=1.55cm and 2.4cm of colmap] (ideal) {Idealbilddomäne\\\texttt{image\_undistorter}\\Maskenwarp mit identischer Abbildung};
      \node[box, below right=1.55cm and 2.4cm of colmap] (gcp) {Originalmodell\\GCP-/CloudCompare-Zweig\\Matrix optional und PA-seitig nicht validiert};
      \node[box, below=1.55cm of ideal] (sts) {\textbf{Container C: STS}\\Objekt-ID-Zuordnung\\Objektbezogene Gaussian-Repräsentation};
      \node[decision, below=1.55cm of sts] (meshroute) {Meshroute};
      \node[route, below left=1.55cm and 1.7cm of meshroute] (original) {\textbf{Route A: Original-GS}\\Surface-Sampling, Poisson, Decimation\\kein SuGaR-Coarse/Refinement/UV};
      \node[route, below right=1.55cm and 1.7cm of meshroute] (coarse) {\textbf{SuGaR-Coarse-Vergleich}\\mask-aware Fork, Coarse, optional Refinement/UV};
      \coordinate (mergebus) at ($(original.south)!0.5!(coarse.south) + (0,-0.55)$);
      \node[box, below=0.95cm of mergebus] (post) {\textbf{Container E: DGtal/GDAL}\\lokales Mesh $\rightarrow$ Centerline\\B-Spline $\rightarrow$ lokale CSV/GeoJSON};
      \node[box, below right=1.8cm and 1.8cm of post] (geo) {Optionaler Exportzweig\\4x4-Transformation + Anchor\\UTM/GIS nur mit validierter Matrix};

      \begin{pgfonlayer}{background}
            \node[container, fill=blue!8, draw=blue!75, fit=(input) (sam)] {};
            \node[container, fill=orange!8, draw=orange!75, fit=(colmap)] {};
            \node[container, fill=cyan!8, draw=cyan!75, fit=(ideal) (sts)] {};
            \node[container, fill=red!8, draw=red!75, fit=(gcp)] {};
            \node[container, fill=purple!8, draw=purple!75, fit=(meshroute) (original) (coarse)] {};
            \node[container, fill=teal!8, draw=teal!75, fit=(post) (geo)] {};
      \end{pgfonlayer}

      \draw[arrow] (input) -- (sam);
      \draw[arrow] (sam) -- (colmap);
      \draw[arrow] (colmap) -- (ideal);
      \draw[arrow] (colmap) -- (gcp);
      \draw[arrow] (ideal) -- (sts);
      \draw[arrow] (sts) -- (meshroute);
      \draw[arrow] (meshroute) -- node[above left] {A} (original);
      \draw[arrow] (meshroute) -- node[above right] {Vergleich} (coarse);
      \draw[arrow] (original) |- (post);
      \draw[arrow] (coarse) |- (post);
      \draw[arrow] (post) -- (geo);
      \draw[arrow] (gcp) |- (geo);
\end{tikzpicture}%
}
\caption{Aktualisierter Datenfluss der fünf Container mit Original-GS-
Produktionsroute, SuGaR-Coarse-Vergleich und optionalem
Georeferenzierungszweig.}
\label{fig:pa-pipeline-flow}
\end{figure}

Das Konzept verfolgt zwei gleichrangige Ziele. Erstens soll ein geführter
Einzellauf von der Eingabe bis zum Endprodukt ohne manuelle Neuverkabelung der
Werkzeuge möglich sein. Zweitens müssen teure Stufen kontrolliert wiederholt
und als Matrix ausgeführt werden können. Jeder Übergang besitzt deshalb
prüfbare Eingaben und persistente Ausgaben.

\subsection{Domänentrennung als zentraler Datenvertrag}

Die methodische Kernidee ist die Domänentrennung: GCP und SfM bleiben im
Originalraum, STS und SuGaR arbeiten in einer idealisierten, konsistenten
Bilddomäne. Dadurch wird vermieden, dass eine Bildverzeichnung nur in den
Kameraparametern, nicht aber in den Masken berücksichtigt wird.

Bei \texttt{SIMPLE\_RADIAL} und \texttt{OPENCV} wird zuerst ein originales
COLMAP-Modell auf den Rohframes geschätzt. Der COLMAP-Undistorter
(\texttt{image\_undistorter}) erzeugt daraus eine separate ideale Szene, indem
er die verzerrten Rohbilder mithilfe des bereits geschätzten Kameramodells
resampled: Jedes Pixel wird an seine radiale Position im verzeichnungsfreien
PINHOLE-Bild neu abgetastet; es wird dabei kein neues Kameramodell geschätzt,
sondern das gelöste Modell angewandt. Die Masken werden mit der identischen
Kameramodellabbildung transformiert. Bei \texttt{PINHOLE} wird dieselbe
Schnittstelle als Pass-through verwendet. Somit sehen STS, Renderer,
SuGaR-Route und Metrikberechnung stets zueinander passende Pixel- und
Kamerakoordinaten.

Diese vollständige Domänenkette ist im automatisierten Matrixpfad, im
restaurierten Smoke-/Replaypfad und im normalen Inline-Vollauf nachgewiesen.
Nach COLMAP rufen Inline-Vollauf und Matrixrunner denselben Maskenwarp mit
identischer Abbildungsvorschrift auf, prüfen die Abdeckung der gewarpten
Masken und übergibt \path{data/03_masks_ideal} an STS und SuGaR. Im
Replay-Dispatcher gilt dasselbe für alle Läufe, die COLMAP selbst ausführen
(Start ab \texttt{gcp}, \texttt{sam3} oder \texttt{colmap}); nur Replays ab
\texttt{sts}, \texttt{sugar} oder \texttt{postprocess} setzen einen von
vornherein passenden idealen Maskensatz voraus, weil dort keine frische
COLMAP-Szene erzeugt wird.

\subsection{Semantische Objektkette}

Die SAM-Masken erfüllen mehrere getrennte Funktionen. Sie definieren die
Objektzuordnung in STS, begrenzen die Auswahl der Ziel-Gaussians und stellen die
Objektregion für die Bildmetriken bereit. In der SuGaR-Coarse-Vergleichsroute
begrenzen sie zusätzlich die photometrische Supervision. Diese mehrfache
Nutzung darf nicht mit einer geometrischen Garantie verwechselt werden: Eine
binäre Maske beschreibt Bildpixel, während ein räumlich ausgedehnter Gaussian
und das aus Samples rekonstruierte Mesh zusätzliche Unsicherheiten besitzen.

\subsection{Produktions- und Vergleichsroute}

Für die Projektarbeit wird Variante A als Produktionsroute verwendet. SuGaR-
Coarse und Refinement bleiben als Vergleichsroute erhalten, sind aber nicht
notwendig, um die direkte STS-Gaussian-Geometrie zu vermessen. Das Mesh ist ein
abgeleitetes Produkt; die Gaussian-Wolke bleibt die primäre geometrische
Hypothese.

Route A führt keinen neuen SuGaR-Splat ein. Sie extrahiert die Oberfläche aus
der vorbereiteten Original-STS-Gaussian-Wolke. Die Vergleichsroute optimiert
dieselbe Ausgangslage zunächst mit dem lokalen maskenbewussten SuGaR-Fork und
kann anschließend ein gebundenes Refinement und UV-Baking ausführen. Die
Trennung ermöglicht, Änderungen der Gaussian-Geometrie von Änderungen der
eigentlichen Mesh-Extraktion zu unterscheiden.

\subsection{Ausführungsformen}

Der interaktive Vollauf führt durch Eingabeauswahl und Parameterentscheidungen.
Vor der Parametrisierung fragt er nach einem Lauf-Preset für die
Arbeitsvideoauflösung (720p, qHD oder low) und weist dabei die gemessenen
Gesamtlaufzeiten der Produktionskonfiguration aus (Route A, 5 FPS,
\texttt{OPENCV}: rund 40 Minuten bei 720p, 34 Minuten bei qHD und 24 Minuten
bei low, jeweils inklusive SAM~3.1-, COLMAP-, Render- und Archivierungsanteil);
das Preset dient der bewussten Wahl zwischen Bildauflösung und Rechenaufwand.
Der Autopilot fragt nach seiner Aktivierung nur noch das Lauf-Preset und den
Objekt-Prompt ab, übernimmt für alle übrigen Entscheidungen validierte
Defaults und überspringt den manuellen
CloudCompare-Breakpoint, wenn keine echte Transformationsmatrix bereitsteht.
Das Preset legt auch die SAM3-Verarbeitungsbreite fest (720p$\to$1280,
qHD$\to$960, low$\to$640 Pixel) und damit die Frame-Auflösung der
Segmentierung. Neben dem Vollprofil existiert das Diagnoseprofil
\texttt{colmap-stop}: Es dient der SfM-Prüfung (Registrierung, Punktanzahl,
Reprojektionsfehler) ohne die teuren Trainingsstufen und warpt die Rohmasken
vor dem Stopp, sodass der Zwischenstand mit einem Replay ab \texttt{sts}
fortsetzbar bleibt.
Der Matrixrunner führt deklarierte Kombinationen dateibasiert
(stdin-unabhängig) nacheinander aus, das heißt er liest die Konfiguration aus
der TSV-Datei und berührt dabei nicht den Standardinput, den nachgelagerte
Unterprozesse für eigene Zwecke konsumieren könnten. Er archiviert jeden Lauf
und setzt den Live-Arbeitsbereich zurück. Diese Modi nutzen dieselben
fachlichen Pipelinekomponenten, erfüllen aber verschiedene Aufgaben der
Benutzerführung und Reproduzierbarkeit.

Die Projektarbeit untersucht die CLI-zentrierte Ausführung. Eine vorhandene
Web-UI und ein tatsächliches Serverdeployment bleiben außerhalb des Scopes.
Die Containerisierung und die nichtinteraktive Matrix zeigen eine für einen
späteren Serverbetrieb geeignete Architektur, aber noch keinen validierten
Produktiv-Rollout.

\subsection{Lokaler Endpunkt und vorbereitete Georeferenzierung}

Die Georeferenzierung ist als nachgelagerter Datenvertrag vorgesehen. Die
Centerline wird zunächst lokal erzeugt. Erst eine validierte Transformation aus
GCP-Picking und ein separater Anchor überführen sie in ein UTM-System. Fehlen
sowohl die Transformationsdatei (\texttt{matrix.txt}) als auch der
Ankerpunkt (\texttt{anchor.txt}), wendet die Pipeline automatisch eine reine
Translation auf einen festen Ersatzanker an und kennzeichnet die Ausgaben mit
dem Suffix \texttt{\_fallback\_georeferenced}. Dieser Fallback dient nur der
technischen Pipelineprüfung und nicht als Genauigkeitsnachweis.
